\begin{remark}[2.5.4]
\label{IV.2.5.4-fr}
Las afirmaciones de (2.5.2) para las propiedades (i) y (ii) siguen siendo
válidas si solo se supone que $f$ es \emph{cuasifielmente plano} (2.3.3) y
cuasicompacto. En efecto, la cuestión se reduce (Bourbaki, \emph{loc. cit.}) a
demostrar el
\end{remark}

\begin{lemma}[2.5.4.1]
\label{IV.2.5.4.1-fr}
Sea $\rho:A\to A'$ un homomorfismo de anillos tal que el morfismo
correspondiente $f:\operatorname{Spec}(A')\to\operatorname{Spec}(A)$ sea
sobreyectivo; supongamos que existe un $A'$-módulo $N'$ de tipo finito, plano
sobre $A$ y cuyo soporte es $\operatorname{Spec}(A')$. Entonces, si
$u:P\to Q$ es un homomorfismo de $A$-módulos tal que
$u\otimes 1:P\otimes_A A'\longrightarrow Q\otimes_A A'$
es sobreyectivo, $u$ es sobreyectivo.
\end{lemma}

En efecto, se deduce primero que el homomorfismo
$u\otimes 1_{N'}:(P\otimes_A A')\otimes_{A'}N'
\longrightarrow(Q\otimes_A A')\otimes_{A'}N'$
es sobreyectivo. Sea $\mathfrak q$ un ideal primo de $A'$ y sea
$\mathfrak p=\rho^{-1}(\mathfrak q)$; el homomorfismo correspondiente
$(P\otimes_A N')_{\mathfrak q}\to(Q\otimes_A N')_{\mathfrak q}$ es
sobreyectivo y se escribe
$u_{\mathfrak p}\otimes 1:
P_{\mathfrak p}\otimes_{A_{\mathfrak p}}N'_{\mathfrak q}
\longrightarrow
Q_{\mathfrak p}\otimes_{A_{\mathfrak p}}N'_{\mathfrak q}$
$(\mathbf{0}_{\mathrm I},1.5.4)$.
Por hipótesis se tiene $N'_{\mathfrak q}\ne0$ y $N'_{\mathfrak q}$ es un
$A_{\mathfrak p}$-módulo plano $(\mathbf{0}_{\mathrm I},6.3.1)$; en virtud del lema de
Nakayama, se tiene
$\mathfrak qN'_{\mathfrak q}\ne N'_{\mathfrak q}$ y, \emph{a fortiori},
$\mathfrak pN'_{\mathfrak q}\ne N'_{\mathfrak q}$, por lo que
$N'_{\mathfrak q}$ es un $A_{\mathfrak p}$-módulo fielmente plano
$(\mathbf{0}_{\mathrm I},6.4.1)$. De ello se sigue que $u_{\mathfrak p}$ es
sobreyectivo $(\mathbf{0}_{\mathrm I},6.4.1)$ y, como esto ocurre para todo
$\mathfrak p\in\operatorname{Spec}(A)$, puesto que $f$ es sobreyectivo, se
concluye finalmente que $u$ es sobreyectivo (Bourbaki, \emph{Alg. comm.},
cap.~II, \S~3, núm.~3, teorema~1).

